\documentclass[a4paper, 10pt]{article}
\usepackage{fontspec}  % Font support for XeLaTeX
\usepackage{geometry}
\geometry{top=1in, bottom=1in, left=1in, right=1in}

\begin{document}

\begin{flushright}
  \texttt{(+86) 193-7488-9316} \\
  \texttt{guanyumao@whu.edu.cn} \\
  %\texttt{https://guanyumao.github.io/}
\end{flushright}\vspace{-45pt}

\begin{flushleft}
  \huge \textbf{Guanyu Mao} \quad \large \textit{Undergraduate Student} \\
  \noindent\rule{\textwidth}{0.4pt}
\end{flushleft}

% -------------------- EDUCATION --------------------
\section*{Education}
\noindent \textbf{Wuhan University}, School of Earth and Space Sciences \hfill Wuhan, China \\
B.S. in Geophysics \hfill Sep 2022 - Jun 2026 \\
GPA: 3.87/4.00 (Rank: 2/23); TOEFL: 85/120

% -------------------- RESEARCH INTERESTS --------------------
\section*{Research Interests}
\begin{itemize}
    \item \textbf{Seismology}: Earthquake rupture processes, dynamic rupture modeling, seismic source inversion, amd induced earthquake.
    \item \textbf{Tectonics}: Fault systems, crustal deformation, plate boundary processes, and volcano-tectonic interactions.
    \item \textbf{Geodynamics}: Stress and strain evolution in the lithosphere, crust-mantle interaction, and deformation mechanisms.
\end{itemize}

% -------------------- RESEARCH EXPERIENCE --------------------
\section*{Research Experience}
\noindent \textbf{Undergraduate Research} \hfill Jan 2023 - Present \\
Supervisor: Prof. Guoyan Jiang, Co-supervised by Postdoc. Yingwen Zhao, Wuhan University
\begin{itemize}
    \item Utilized GAMMA software to process D-InSAR data (Sentinel-1) and extract coseismic surface deformation associated with the 2023 Xinwen Ms 4.9 earthquake sequence.  
    \item Conducted inversion of fault geometry and slip distribution based on the derived surface deformation models.  
    \item Analyzed earthquake catalogs to investigate spatiotemporal patterns of aftershocks.
    \item Researched corelations between aftershock activity and seismogenic fault.
    \item Applied the Epidemic-Type Aftershock Sequence (ETAS) model to assess earthquake triggering behavior within the sequence.  
    \item Performed GNSS velocity field analysis to study regional crustal strain distribution.
\end{itemize}

% -------------------- TECHNICAL SKILLS --------------------
\section*{Technical Skills}
\noindent \textbf{Programming}: Python, MATLAB, C, C++, C\#, Shell \\
\textbf{Tools and Software}: GMT, Git, LATEX, Markdown, HTML, Gamma\\
\textbf{Languages}: Chinese (native), English (fluent)

% -------------------- PUBLICATIONS --------------------
%\section*{Publications}
%\textit{Manuscript in preparation} on the 2023 Xinwen earthquake sequence, including InSAR deformation modeling, fault inversion, and aftershock analysis.

\end{document}
